\lecture{21}{May 22.}{}
\section{Edge Coloring}
\begin{definition}
    Given a graph \(G = (V, E)\), a proper \(k\)-edge-coloring is a map \(c: E \to [k]\) such that any pair of incident edges \((e, f \in E, e \cap f \neq \varnothing)\) receive disitnct colors.    
\end{definition}

\begin{definition}
    The edge-chromatic number (chromatic index) \(\chi ^{\prime} (G)\) is the minimum \(k\) s.t. \(G\) has a proper \(k\)-edge-coloring.    
\end{definition}

\begin{observation}
    In a proper edge-coloring, each color class is a matching, so proper edge-coloring is a decomposition into matchings.
\end{observation}

\begin{proposition}
    For any graph \(G\), \(\chi ^{\prime} (G) \ge \Delta (G)\).  
\end{proposition}

\begin{proof}
    The edges incident to a vertex of maximum degree must all receive distinct colors. 
\end{proof}

\begin{theorem}[Konig]
    If \(G\) is bipartite, then \(\chi ^{\prime} (G) = \Delta (G)\).  
\end{theorem}

\begin{proof}
    We have the following cases:

    \textbf{Case 1:} \(G\) is \(\Delta \)-regular, i.e. every vertex has degree \(\Delta\). 

    \begin{claim}
        \(G\) has a perfect matching. 
    \end{claim}    

    \begin{explanation}
        Applying Hall's theorem: \(\forall S \subseteq X\),
        \[
            \vert S \vert \Delta = e(S, N(S)) \le \vert N(S) \vert \Delta.  
        \] 
        Hence, \(G\) has a perfect matching. 
    \end{explanation}

    Hence, we can make this one color, remove, and repeat. 

    \textbf{Case 2:} \(G\) is not \(\Delta \)-regular. 

    We can embed\footnote{HW} \(G \subseteq G^{\prime}\), a \(\Delta \)-regular graph. (draw some edges on \(G\) to make it cbecome \(G^{\prime}\))    
\end{proof}

\paragraph{Line graphs}
Given \(G=(V, E)\), the line graph \(L(G)\) has vertices \(V(L(G)) = E\) and edges 
\[
    E(L(G)) = \left\{ \left\{ e, f \right\}: e \cap f \neq \varnothing  \right\}. 
\]   

\begin{observation}
    \(\chi^{\prime} (G) = \chi (L(G))\), i.e. proper edge-coloring of \(G\) corresponds to a proper (vertex)-coloring of \(L(G)\).     
\end{observation}

\begin{observation}
    Every vertex in \(G\) gives a clique in \(L(G)\), and \(K_3\) in \(G\) gives \(K_3\) in \(L(G)\). Hence, \(\omega (L(G)) = \Delta (G)\) if \(\Delta (G) \ge 3\).          
\end{observation}

We therefore have 
\begin{align*}
    \chi ^{\prime} (G) &= \chi (L(G)) \le \Delta (L(G)) + 1 \le (2 \Delta (G) - 2) + 1 = 2 \Delta (G) - 1.
\end{align*}
Note that 
\[
    \Delta (L(G)) \le 2 \Delta (G) - 2
\]
since \(\Delta (L(G))\) is the maximum number of edges in \(G\) incident to the endpoint of an edge. 

Hence,
\[
    \Delta (G) \le \chi ^{\prime} (G) \le 2 \Delta (G) - 1.
\]

\begin{theorem}[Vizing]
    For any \(G\), 
    \[
      \Delta (G) \le \chi ^{\prime} (G) \le \Delta (G) + 1.  
    \] 
\end{theorem}

\begin{definition}
    A graph \(G\) is Type 1 if \(\chi ^{\prime} (G) = \Delta (G)\), and Type 2 if \(\chi ^{\prime} (G) = \Delta (G) + 1\).   
\end{definition}

\begin{eg}
    All bipartite graphs are Type 1. Even cliques are Type 1. Odd cycles and odd cliques are Type 2.
\end{eg}

\begin{proof}[proof of Vizing]
    Let \(v \in V(G)\), consider \(G^{\prime} = G - v\). By induction, we can \(\Delta (G^{\prime} ) + 1 \le (\Delta (G) + 1)\)-edge-color \(G^{\prime}\). 
    
    We need to color the edges incident to \(v\).      

    \begin{lemma}
        Let \(G\) be a graph, \(v \in V(G)\), \(k \in \mathbb{N}\). Then, if:
        \begin{itemize}
            \item \(d(v) \le k\). 
            \item for each neighbor \(u\) of \(v\), \(d(u) \le k\). 
            \item there is at most one neighbor \(u\) of \(v\) with \(d(u) = k\).       
        \end{itemize}   
        Then, \(G\) is \(k\)-edge-colorable iff \(G - v\) is \(k\)-edge-colorable.    
    \end{lemma}

    To deduce the theorem, we apply the lemma with \(k = \Delta (G) + 1\).  
\end{proof}